% !Mode:: "TeX:UTF-8"
% VL 3, gebaut aus part-grundlagen/04-vortrainierte-modelle.md und 05-tensoren.md
% (Dialogmuster und Zeiten: transcripts/vl03-pytorch-tensoren-2025/notiz.md).
% !Mode:: "TeX:UTF-8"
\newcommand{\thepage}{} % to silence a hyperref warning about pdfpagelabels
\RequirePackage[]{silence}
\WarningsOff[hyperref]
\PassOptionsToPackage{log-declarations=false}{xparse}
%\RequirePackage[l2tabu, orthodox]{nag}
\RequirePackage{luatex85} % to make \pdf...= commands work with LuaTeX
\documentclass[aspectratio=169,t]{beamer}
\usepackage{ifthen}
\usepackage[T1]{fontenc}
\usepackage{iftex}
\ifLuaTeX
 \usepackage{fontspec}
 \setmainfont[Ligatures=TeX]{Source Sans Pro}
\else
 \usepackage[utf8]{inputenc}
\fi
\usepackage{german}
\usepackage{newunicodechar}
\newunicodechar{€}{\texteuro}
\newunicodechar{ß}{\ss}
\renewcommand{\SS}{\iffontchar\font"1E9E \symbol{"1E9E}\else SS\fi}
\usepackage{ellipsis}
\usepackage{microtype}
\usepackage{listings}
\usepackage{sourcecodepro}
\lstset{language=Python, basicstyle=\ttfamily}
\usepackage{amsfonts}
\usepackage{amssymb}
\usepackage{amsthm}
\usepackage{fancyhdr}
\usepackage{graphicx}

\newtheorem{proposition}[theorem]{Proposition}

\DeclareMathOperator*{\argmax}{argmax}
\DeclareMathOperator{\supp}{supp}
\newcommand{\E}{\ensuremath{\mathbb{E}}}
\newcommand{\V}{\ensuremath{\mathbb{V}}}

\setbeamercolor{emph}{fg=purple}
\renewcommand<>{\emph}[1]{%
  {\usebeamercolor[fg]{emph}\only#2{\itshape}#1}%
}



\hypersetup{draft=false,colorlinks=false,pdfpagemode=UseOutlines,plainpages=false,
            breaklinks=true,bookmarksnumbered=true,pdfborder=0 0 0,
            pdftitle={Einführung in Deep Learning WiSe 2026/27 HHU},
            pdfauthor={Konrad Voelkel},
            pdfsubject={data science}}

%\documentclass[aspectratio=169,t]{beamer}
\usetheme{hhubeamer} % put the hhu-beamer-theme in /usr/share/texmf/tex/latex/beamer/themes/

% -------------------- Code listings config --------------------
% Define custom colors for syntax highlighting
\definecolor{codegreen}{rgb}{0,0.6,0}
\definecolor{codegray}{rgb}{0.5,0.5,0.5}
\definecolor{codepurple}{rgb}{0.58,0,0.82}
\definecolor{backcolour}{rgb}{0.96,0.96,0.94} % Very light gray/beige
% Configure the listings package
\lstset{
    language=Python,
    backgroundcolor=\color{backcolour},   
    commentstyle=\color{codegreen}\itshape,
    keywordstyle=\color{blue}\bfseries,
    numberstyle=\tiny\color{codegray},
    stringstyle=\color{codepurple},
    basicstyle=\ttfamily\footnotesize, % Use a monospaced font
    breakatwhitespace=false,         
    breaklines=true,      
             % Enable automatic line breaking
    captionpos=b,                      % Caption at the bottom
    keepspaces=true,                 
    numbers=left,                      % Line numbers on the left
    numbersep=10pt,                    % Distance of line numbers from the code
    showspaces=false,                
    showstringspaces=false,            % Don't show weird underscores in strings
    showtabs=false,                  
    tabsize=4,
    frame=single,                      % Add a subtle border around the code
    rulecolor=\color{black!20},        % Border color
  }

\event[]{}
\author[Wintersemester 2026/27]{Dr. Konrad Völkel, Wintersemester 2026/27}

%\title{Data Science, VL 1}
%\date{16.04.2026}


%%% Local Variables:
%%% mode: latex
%%% TeX-master: t
%%% End:


% Fassung mit Dozentennotizen (zweiter Schirm): die nächste Zeile einkommentieren.
% \setbeameroption{show notes on second screen=right}

\graphicspath{{../images/}}

\title[Vortrainierte Modelle und Tensoren (VL 3)]{Vortrainierte Modelle und Tensoren\\\large Einführung in Deep Learning, VL 3}
\date{28.10.2026}

\begin{document}

\begin{frame}
  \titlepage

  \note{\textbf{Ablauf, 90 Minuten} (2025: 55 Min. Folien, 14 Min. Code, 21 Min. Tensor-Ausflug
    mit Tafel):\\[2pt]
    \begin{tabular}{@{}rll@{}}
      0--4   & Folie 2       & Rückblick Übungsblatt, heute zwei Themen\\
      4--12  & Folien 3--4   & Hugging Face: was es ist, wie man sucht\\
      12--30 & Folien 5--7   & \textbf{Task-Rundgang} (Browser), Vokabular\\
      30--42 & Folien 8--9   & WMT19-Beispiel, das Muster Tokenizer--Modell--Tokenizer\\
      42--50 & Folien 10--11 & GPU, Installation, Colab\\
      50--58 & Folien 12--14 & Tensor $\approx$ Array, und was er mehr hat\\
      58--75 & Folien 15--17 & \textbf{Ausflug Tensorprodukt (Tafel)}, Reshape-Quiz\\
      75--85 & Folien 18--19 & Das Batch-Bild, Shape-Fehler\\
      85--90 & Folie 20      & Übungsblatt\\
    \end{tabular}\\[4pt]
    Der Ausflug ist der anspruchsvollste Teil der Sitzung und lebt von der Tafel.
    2025 war der Stoff nach gut einer Stunde durch und der Rest wurde Vorgriff --
    lieber den Ausflug langsamer machen als vorzugreifen.}
\end{frame}

\begin{frame}
  \frametitle{Wo wir stehen}

  \begin{center}
    \Large
    Bisher: ein Neuron, das wir selbst trainieren.\\[14pt]
    Heute: Modelle, die \emph{andere} trainiert haben.
  \end{center}

  \note{Das Perzeptron-Blatt ist abgegeben, das nächste kommt raus.
    Heute zwei Themen, die wenig miteinander zu tun scheinen und es doch tun: Wer fertige
    Modelle benutzt, bekommt Tensoren zurück und muss sie lesen können.}
\end{frame}

% ------------------------------------------------------------------ Hugging Face
\begin{frame}
  \frametitle{Hugging Face}

  \begin{center}
    \Large
    \href{https://huggingface.co/models}{\color{eidlbrand}huggingface.co/models}\\[6pt]
    \href{https://huggingface.co/tasks}{\color{eidlbrand}huggingface.co/tasks}
  \end{center}

  \bigskip
  \begin{itemize}
  \item Firma aus New York, 2016 gegründet
  \item verdient Geld mit gehosteter Inferenz
  \item lockt mit Infrastruktur: \texttt{transformers}, \texttt{diffusers}, Hosting, Doku
  \item September 2026: Nvidia kauft Hugging Face für fast 13 Mrd. \$ (ca. 11,6 Mrd. €)
  \end{itemize}

  \note{\textbf{Beide Adressen sind anklickbar} -- hier in den Browser wechseln und ein paar
    Minuten herumklicken, der Rundgang ab Folie 5 lebt davon.
    Der Hub ist ein kommerzielles Angebot, aber die Modelle und Libraries darauf sind
    größtenteils offen -- und für uns die bequemste Art, an Gewichte zu kommen.
    In der Finanzierungsrunde vom August 2023 war die Firma 4,5 Mrd. \$ wert; im September 2026
    hat Nvidia den Kauf für 12,93 Mrd. \$ vereinbart -- der Abschluss wird für die erste
    Hälfte 2027 erwartet, vorbehaltlich der Freigaben. Also knapp das Dreifache in drei Jahren,
    und die Plattform mit den offenen Gewichten gehört künftig demselben Konzern wie die
    Hardware, auf der sie gerechnet werden. Das ist die Einordnung, nicht nur die Zahl.
    2025 hat hier der Vergleich Likes gegen Downloads gut funktioniert: Likes sagen, was
    die Community feiert, Downloads, was tatsächlich benutzt wird.}
\end{frame}

\begin{frame}
  \frametitle{Wonach sortiert man Modelle?}

  \begin{center}
    \Large Nach \emph{Task}.
  \end{center}

  \bigskip
  \begin{center}
    Computer Vision \quad NLP \quad Audio\\[6pt] Multimodal \quad Tabular \quad Reinforcement Learning
  \end{center}

  \note{Jetzt in den Browser: huggingface.co/tasks. Der Rundgang dient nicht dem
    Verständnis der Modelle, sondern dem \emph{Vokabular} -- das ist die Sprache, in der
    Modellkarten, Papers und Stellenanzeigen reden.
    Audio, Tabular und Reinforcement Learning klammern wir im Kurs aus, das hier ansagen.}
\end{frame}

\begin{frame}
  \frametitle{Computer Vision}

  % description statt itemize: Beamer schneidet lange \item[...]-Labels links ab.
  \begin{description}[Image Classification]\setlength{\itemsep}{4pt}
  \item[Depth Estimation] pro Pixel eine $z$-Koordinate \hfill {\small DPT}
  \item[Image Classification] was ist drauf? \hfill {\small ResNet}
  \item[Object Detection] was ist wo? \hfill {\small YOLOS}
  \item[Segmentation] Umrisse pro Objekt \hfill {\small Segformer}
  \item[Image-to-Image] Bild rein, Bild raus \hfill {\small ColorizeNet}
  \item[Zero-Shot] Klassen als Text \hfill {\small CLIP-ViT}
  \end{description}

  \note{Im Browser jeweils ein Beispiel live zeigen -- die Task-Seiten haben Widgets.
    Vokabel erklären: \emph{Zero-Shot} heißt ohne gelabelte Daten für die neue Aufgabe,
    \emph{One-Shot} mit einem Sample, \emph{Few-Shot} mit einigen wenigen.
    2025 kam die Frage, ob Object Detection ein Spezialfall von Image Classification sei.
    Antwort: unterschiedliche Aufgabenstellung, aber es gibt Multi-Task-Modelle.}
\end{frame}

\begin{frame}
  \frametitle{NLP}

  \begin{description}[Token Classification]\setlength{\itemsep}{6pt}
  \item[Text Classification] Genre, Sentiment \hfill {\small DistilBERT}
  \item[Token Classification] NER, PoS-Tagging \hfill {\small flair ner-german}
  \item[Summarization] Zusammenfassung \hfill {\small BART}
  \item[Fill-Mask] Lückentext \hfill {\small DistilBERT}
  \end{description}

  \note{Fill-Mask ist die interessanteste Zeile: Das ist genau die Aufgabe, mit der
    BERT-artige Modelle \emph{vortrainiert} werden -- Lücke rein, vorhersagen lassen.
    Diese Idee kommt bei word2vec (VL 9) wieder.
    DistilBERT ist ein destilliertes Modell; Destillation ist VL 15.}
\end{frame}

\begin{frame}
  \frametitle{Multimodal}

  \begin{description}[Document Question Answering]\setlength{\itemsep}{6pt}
  \item[Image-to-Text] Bildbeschreibung \hfill {\small BLIP}
  \item[Visual Question Answering] Frage zum Bild \hfill {\small ViLT}
  \item[Document Question Answering] Frage ans Formular \hfill {\small LayoutLM}
  \item[Feature Extraction] Vektor statt Antwort \hfill {\small CLIP}
  \end{description}

  \note{Multimodal heißt: mehr als eine Sorte Eingabe. Das ist die Schublade, in der die
    Modelle liegen, die Bild und Text zugleich sehen -- CLIP kam auf der Vision-Folie schon
    als Zero-Shot-Klassifikator vor, hier steht dasselbe Modell als Merkmalsextraktor.
    \emph{Feature Extraction} ist die Zeile, auf die es ankommt: Das Modell gibt keine
    Antwort, sondern einen Vektor -- Repräsentationen, der rote Faden ab VL 9.
    Audio, Tabular und Reinforcement Learning lassen wir aus, wie auf der Übersicht angesagt.}
\end{frame}

% ------------------------------------------------------------------ Praxis
\begin{frame}[fragile]
  \frametitle{Praktisch: Übersetzen}

  \begin{lstlisting}
from transformers import FSMTForConditionalGeneration, FSMTTokenizer

mname = "facebook/wmt19-en-de"
tokenizer = FSMTTokenizer.from_pretrained(mname)
model = FSMTForConditionalGeneration.from_pretrained(mname)

input = "I am not sure if deep learning is really for me."
input_ids = tokenizer.encode(input, return_tensors="pt")
outputs = model.generate(input_ids)
print(tokenizer.decode(outputs[0], skip_special_tokens=True))
\end{lstlisting}

  \begin{center}
    \footnotesize\texttt{Ich bin mir nicht sicher, ob Deep Learning wirklich etwas für mich ist.}
  \end{center}

  \begin{lstlisting}
uv run --with 'transformers[torch]' --with sacremoses uebersetzen.py
\end{lstlisting}

  \note{Die Ausgabe steht auf der Folie und ist am 2026-10-28 nachgefahren worden, der Satz
    stimmt wörtlich.
    \texttt{sacremoses} gehört dazu: Der FSMT-Tokenizer bricht ohne das Paket mit einem
    \texttt{ImportError} ab, und \texttt{transformers[torch]} zieht es nicht mit -- das Paket
    steht nur in den Extras \texttt{testing} und \texttt{dev}, also muss man es nennen.
    Nachgefahren am 2026-10-28 mit transformers 5.21.0, torch 2.14.1, huggingface\_hub 2.3.0.
    WMT19 ist nach der Konferenz benannt (Workshop on Machine Translation).
    \texttt{from\_pretrained} lädt beim ersten Aufruf die Gewichte aus dem Hub und cached
    sie lokal -- das erste Mal dauert.
    Die zweite Zeile ist der kürzeste Weg zum Ausprobieren: \texttt{uv run -{}-with} legt die
    Umgebung für diesen einen Lauf an und wirft sie danach weg -- kein venv, kein
    \texttt{pip install}. Der dauerhafte Weg steht auf Folie 11.}
\end{frame}

\begin{frame}
  \frametitle{Das Muster hinter jedem Hub-Modell}

  \begin{center}
    \setlength{\tabcolsep}{4pt}
    \begin{tabular}{@{}ccccccc@{}}
      {\LARGE Daten} & {\LARGE $\to$}   & {\LARGE Zahlen} & {\LARGE $\to$} & {\LARGE Zahlen} & {\LARGE $\to$}   & {\LARGE Daten}\\[4pt]
                    & \small Tokenizer &                & \small Modell &                & \small Tokenizer &
    \end{tabular}
  \end{center}

  \medskip
  \begin{center}
    \small Bei Bild und Audio heißt der erste Schritt nicht Tokenizer, sondern \emph{Prozessor}.
  \end{center}

  \note{Das Modell rechnet nur mit Zahlen; der Tokenizer ist die Übersetzung in beide
    Richtungen. Das gilt für jede Eingabe, nicht nur für Text -- deshalb steht oben
    \emph{Daten} und nicht \emph{Text}: Ein Bild wird zu Pixelwerten, Audio zu einem
    Spektrogramm, und in \texttt{transformers} heißt das Gegenstück zum Tokenizer dann
    \texttt{AutoProcessor} oder \texttt{FeatureExtractor}. Gemeinsam ist allen drei Fällen:
    kodieren, rechnen, dekodieren.
    \texttt{return\_tensors="pt"} heißt: gib mir PyTorch-Tensoren -- damit
    sind wir beim zweiten Thema.
    Kurzer Ausblick, falls die Frage kommt: Warum nicht Wörter als Atome? Weil das
    Vokabular explodiert und unbekannte Wörter keine Repräsentation hätten. Details in VL 10.}
\end{frame}

\begin{frame}[fragile]
  \frametitle{Auf der GPU rechnen}

  \begin{lstlisting}
import torch

device = "cuda" if torch.cuda.is_available() else "cpu"
model.to(device)
input_ids = input_ids.to(device)
  \end{lstlisting}

  \bigskip
  \begin{center}
    Modell \emph{und} Daten müssen auf dasselbe Gerät.
  \end{center}

  \note{Der häufigste Anfängerfehler ist, nur eines von beiden zu verschieben -- die
    Fehlermeldung dazu ist aber eindeutig.
    Auf AMD heißt es ROCm, der Python-Befehl bleibt derselbe; auf macOS gibt es MPS.
    Wir programmieren kein CUDA, nur den Wrapper PyTorch -- das 2025 ausdrücklich
    versprochen.}
\end{frame}

\begin{frame}[fragile]
  \frametitle{Installation}

  \textbf{Ohne uv}
  \begin{lstlisting}
python -m venv .env
source .env/bin/activate
pip install 'transformers[torch]'
\end{lstlisting}

  \smallskip
  \textbf{Mit uv}
  \begin{lstlisting}
uv init
uv add 'transformers[torch]'
uv sync
\end{lstlisting}

  \smallskip
  \begin{center}
    Ohne GPU? \quad Google Colab.
  \end{center}

  \note{Handzeichen einsammeln: Wer bekommt lokal \texttt{import torch} zum Laufen?
    2025 waren es zwei Personen -- deshalb Colab ausdrücklich als gleichwertigen Weg
    anbieten, es gibt dort begrenzte GPU-Zeit im Browser.
    Der Unterschied: Links legt man die Umgebung selbst an und muss sie vor jedem Start
    aktivieren. Rechts schreibt \texttt{uv add} das Paket in die \texttt{pyproject.toml} und
    hält es in \texttt{uv.lock} auf einer festen Fassung fest; \texttt{uv sync} baut die
    Umgebung daraus -- deshalb bekommt ein zweiter Rechner mit denselben zwei Dateien
    dieselbe Umgebung, und \texttt{uv run} aktiviert nichts von Hand.
    Das \texttt{uv run -{}-with} von Folie 8 ist der dritte Fall: kein Projekt, nur einmal
    ausprobieren.
    Für CUDA: \texttt{nvidia-smi} zeigt Treiber und maximale CUDA-Version, dann den
    passenden torch-Build installieren; \texttt{torch.cuda.is\_available()} muss True sagen.}
\end{frame}

% ------------------------------------------------------------------ Tensoren
\begin{frame}
  \frametitle{Vier Fragen}

  \begin{itemize}\setlength{\itemsep}{8pt}
  \item Was ist ein Array?
  \item Was ist ein Tensor?
  \item Was ist ein Array in NumPy?
  \item Was ist ein Tensor in Torch?
  \end{itemize}

  \note{Array: Daten gleichen Typs, fortlaufend im Speicher.
    Tensor, je nachdem wen man fragt: Physiker -- etwas, das sich tensoriell
    transformiert; Mathematiker -- Element eines Tensorproduktraums; Informatiker --
    ein fancy Wort für Array.
    Torch-Tensor: Daten gleichen \texttt{dtype}, fortlaufend, potentiell im GPU-Speicher,
    optimiert für Gradientenabstieg (Attribut \texttt{grad}).}
\end{frame}

\begin{frame}
  \frametitle{Tensor $\approx$ Array}

  \begin{center}
    \Large \texttt{torch.from\_numpy(a)}
  \end{center}

  \bigskip
  \begin{center}
    teilt sich den Speicher mit dem NumPy-Array.
  \end{center}

  \note{Kein Kopieren, kein zusätzlicher Speicher -- eine Änderung am einen ist eine
    Änderung am anderen. Das ist Segen und Falle zugleich.
    Empfehlung: das PyTorch-Tutorial „Tensors“ einmal durchklicken, steht im Skript verlinkt.}
\end{frame}

\begin{frame}
  \frametitle{Was ein Torch-Tensor mehr hat}

  \begin{center}
    \begin{tabular}{@{}lll@{}}
                            & NumPy            & PyTorch\\[4pt]
      \texttt{dtype}, \texttt{shape} & ja        & ja\\[4pt]
      \texttt{device}       & immer \texttt{cpu} & \texttt{cpu}, \texttt{cuda}, \texttt{mps}\\[4pt]
      \texttt{requires\_grad} & --             & wird nach mir abgeleitet?\\[4pt]
      \texttt{grad}         & --               & hier landet die Ableitung\\[4pt]
      \texttt{grad\_fn}     & --               & woraus ich entstanden bin\\
    \end{tabular}
  \end{center}

  \bigskip
  \begin{center}
    Dieselben Daten, mehr Buchhaltung: Der Tensor merkt sich seine Herkunft.
  \end{center}

  \note{Das ist die Antwort auf die vierte der vier Fragen, und sie ist der Grund, warum es
    Torch-Tensoren überhaupt gibt -- die Daten könnte NumPy auch halten.
    Vorführen, wenn Zeit ist: \texttt{w = torch.tensor([2.0], requires\_grad=True)},
    \texttt{y = (w * 3).sum()}. Dann ist \texttt{w.grad} noch \texttt{None} und
    \texttt{y.grad\_fn} ein \texttt{SumBackward0}; nach \texttt{y.backward()} steht in
    \texttt{w.grad} der Wert \texttt{tensor([3.])}. Am 2026-10-28 mit torch 2.14.1 so
    nachgefahren.
    Die drei Spalten unten sind der rote Faden des Semesters: \texttt{grad} ist der
    Gradientenabstieg (VL 5), \texttt{grad\_fn} der Berechnungsgraph hinter Backpropagation
    (VL 6) und Autograd (VL 12). Hier nur zeigen, dass es die Felder gibt -- nicht erklären,
    wie sie gefüllt werden.
    \texttt{device} steht übrigens auch in NumPy, seit es dem Array-API-Standard folgt, dort
    aber immer auf \texttt{cpu}. Die Wahl des Geräts ist der zweite Unterschied, siehe Folie 10.}
\end{frame}

\begin{frame}
  \frametitle{Ausflug: Tensorprodukt}

  \begin{center}
    \Large
    $\{$bilineare Abb. $V \times W \to U\}
    \;\longleftrightarrow\;
    \{$lineare Abb. $V \otimes W \to U\}$
  \end{center}

  \bigskip
  \begin{center}
    $\dim(V \otimes W) = n \cdot m$, \qquad $\mathbb{R}^n \otimes \mathbb{R}^m \cong \mathbb{R}^{nm}$
  \end{center}

  \note{\textbf{Tafelphase.} Ankündigen, dass jetzt ein Ausflug kommt, den man überspringen
    darf -- und wofür er gut ist: um zu verstehen, was \texttt{shape} und \texttt{reshape}
    eigentlich tun.
    An der Tafel: Skalare im Tensorprodukt hin- und herschieben, damit die Bilinearität
    anschaulich wird. Basen $v_i \otimes w_j$ zählen.
    Wer nicht mehr folgt: Gedanken auf die Übungsaufgabe schreiben, wir klären es dort.}
\end{frame}

\begin{frame}
  \frametitle{Ausflug: Dualraum}

  \begin{center}
    \Large
    $V^\vee = \operatorname{Hom}(V, \mathbb{R})$\\[14pt]
    \normalsize
    Spaltenvektor $v$ \quad$\longleftrightarrow$\quad Zeilenvektor $v^T$\\[8pt]
    $v^T w = v \cdot w$
  \end{center}

  \note{An der Tafel: der Kovektor als „liegender Vektor“, geschrieben als
    $a x + b y + c z$ -- eine Maschine, die einen Vektor frisst und eine Zahl ausspuckt.
    Kovariant und kontravariant nur benennen, nicht vertiefen; in Einstein-Notation stehen
    die Indizes oben oder unten. Die Physiker im Deep Learning bringen diese Sprache mit.}
\end{frame}

\begin{frame}
  \frametitle{Quizfrage}

  \begin{center}
    \Huge Was macht \texttt{reshape}?
  \end{center}

  \note{Sammeln, nicht sofort auflösen. 2025 kamen Transponieren und lineare Abbildung --
    beides aufgreifen und abgrenzen.
    Auflösung am Beispiel des Kapitels: $a \in \mathbb{R}^4 \otimes \mathbb{R}^{10} \otimes
    \mathbb{R}^2$, das sind achtzig Zahlen. Wegen
    $\mathbb{R}^4 \otimes \mathbb{R}^{10} \cong \mathbb{R}^{40}$ ist
    \texttt{a.reshape(-1, 2)} dasselbe wie \texttt{a.reshape(40, 2)} -- derselbe Tensor,
    nur mit anderer Klammerung der Faktoren.
    \emph{Reshape ändert keine Daten, nur die Interpretation des Speichers.}
    An der Tafel die achtzig Zahlen anders umbrechen.}
\end{frame}

\begin{frame}
  \frametitle{Das Bild, das man sich merken sollte}

  \begin{center}
    \Large
    \texttt{(784,)} \quad \emph{ein} Datenpunkt, 784 Features\\[8pt]
    \texttt{(16, 784)} \quad ein Batch aus 16 Datenpunkten\\[8pt]
    \texttt{(16, 28, 28)} \quad dieselben 16 als Bilder\\[8pt]
    \texttt{(16, 3, 28, 28)} \quad dieselben mit Farbkanälen
  \end{center}

  \medskip
  \begin{center}
    Die \emph{erste} Achse ist der Batch.\\[6pt]
    \small 32 Bilder bei \texttt{batch\_size=16}: \emph{zwei} Tensoren \texttt{(16, 28, 28)}.\\[4pt]
    \small PIL, matplotlib: \texttt{(H, W, C)} \qquad PyTorch: \texttt{(N, C, H, W)}
  \end{center}

  \note{Das ist das mentale Bild: gestapelte Matrizen, und der vorderste Index sagt,
    das wievielte Bild im \emph{Batch} es ist.
    \textbf{Der Punkt, an dem es klemmt:} Die erste Zahl ist nicht die Größe des Datensatzes,
    sondern wie viele Beispiele \emph{gerade zusammen} durch das Netz gehen. 32 Bilder bei
    \texttt{batch\_size=16} sind also keine \texttt{(32, …)}, sondern zwei Tensoren
    \texttt{(16, 28, 28)}, die der \texttt{DataLoader} nacheinander ausgibt. Bei 30 Bildern
    wäre der zweite \texttt{(14, 28, 28)} -- der letzte Batch darf kleiner sein, es sei
    denn, man setzt \texttt{drop\_last=True}. Den ganzen Datensatz als einen Tensor gibt es
    auch, man füttert das Modell aber nicht damit.
    \textbf{Der Kanalvergleich unten:} Ein Bild aus PIL, matplotlib oder \texttt{skimage}
    liegt als $(H, W, C)$ vor, Kanäle hinten; PyTorch erwartet $(N, C, H, W)$, Kanäle vorn.
    Deshalb steht in jedem zweiten Bildskript ein \texttt{permute(2, 0, 1)}. Dass Zeilen die
    Datenpunkte sind, ist dagegen in beiden Welten gleich (NumPy, scikit-learn:
    \texttt{(n\_samples, n\_features)}) -- daran liegt der Unterschied also nicht.
    \texttt{nn.Linear} und fast alle Layer wirken auf die \emph{letzte} Achse und lassen
    den Batch in Ruhe. Das „fast“ sind die rekurrenten Schichten: \texttt{nn.RNN},
    \texttt{nn.LSTM}, \texttt{nn.GRU} und \texttt{nn.Transformer} erwarten voreingestellt
    \texttt{(seq, batch, feature)}, legen den Batch also auf die \emph{zweite} Achse;
    \texttt{batch\_first=True} stellt die Regel wieder her. Falls gefragt wird, warum hier
    \texttt{(32, 3, 28, 28)} und nicht \texttt{(32, 28, 28, 3)}: Die Faltungsschichten (VL 11)
    verlangen die Kanalachse vorn, $(N, C, H, W)$.}
\end{frame}

\begin{frame}
  \frametitle{Warum das wichtig ist}

  \begin{center}
    \Large Der Shape sagt, \emph{wie viele} Zahlen in jede Richtung gehen --\\[8pt]
    nicht, \emph{was} die Richtung bedeutet.
  \end{center}

  \note{Das ist die Take-Home-Message der Sitzung. Die Anekdote erzählen: einmal bei einem
    \emph{quadratischen} Bild die beiden Achsen vertauscht -- kein Fehler, keine Warnung,
    nur wochenlang sinnloser Output.
    Wer einen Tensor vor sich hat, sollte als Erstes sagen können: Welche Achse ist der
    Batch, was bedeuten die übrigen?}
\end{frame}

\begin{frame}
  \frametitle{Übungsblatt}

  \begin{itemize}\setlength{\itemsep}{8pt}
  \item Ableitungen wiederholen -- und an einer Stelle verallgemeinern
  \item Die Betragsfunktion ist banal, aber wichtig
  \item Die Integrale dürfen Sie beim ersten Lesen ignorieren
  \end{itemize}

  \bigskip
  \begin{center}
    Nächste Woche: logistische Regression.
  \end{center}

  \note{Zur Betragsfunktion: Man soll verstehen, dass die an der Null umdefinierte
    Vorzeichenfunktion ebenfalls eine schwache Ableitung ist -- daher kommt später ReLU
    davon, obwohl sie in 0 nicht differenzierbar ist.
    Warum man die Integrale ignorieren darf, besprechen wir in der Übung.}
\end{frame}

\end{document}

%%% Local Variables:
%%% mode: latex
%%% TeX-master: t
%%% End:
